\documentclass[a4paper,12pt]{ctexart}
\usepackage{xcolor,graphicx,amsmath}
\usepackage[top=1.8cm, bottom=1.8cm, left=1.8cm, right=1.8cm]{geometry}
\usepackage{array,hyperref}
\hypersetup{colorlinks=true,linkcolor=blue!70!black}
\linespread{1.35}
\definecolor{defblue}{RGB}{0,80,160}\definecolor{thinkorange}{RGB}{230,120,20}
\definecolor{funboxbg}{RGB}{245,248,252}\definecolor{practicegreen}{RGB}{40,130,70}
\newenvironment{thinkbox}{\vspace{0.3cm}\begin{center}\begin{minipage}{0.88\textwidth}\colorbox{funboxbg}{\begin{minipage}{0.94\textwidth}\vspace{0.15cm}\textbf{\textcolor{thinkorange}{【 想一想 】}}}{\vspace{0.15cm}\end{minipage}}\end{minipage}\end{center}\vspace{0.3cm}}
\newenvironment{definition}[1]{\vspace{0.2cm}\begin{center}\begin{minipage}{0.88\textwidth}\fbox{\begin{minipage}{0.92\textwidth}\vspace{0.1cm}\textbf{\textcolor{defblue}{【 #1 】}}}{\vspace{0.1cm}\end{minipage}}\end{minipage}\end{center}\vspace{0.2cm}}
\newenvironment{practice}{\vspace{0.3cm}\begin{center}\begin{minipage}{0.88\textwidth}\colorbox{funboxbg}{\begin{minipage}{0.94\textwidth}\vspace{0.15cm}\textbf{\textcolor{practicegreen}{【 练一练 】}}}{\vspace{0.15cm}\end{minipage}}\end{minipage}\end{center}\vspace{0.2cm}}
\newenvironment{figplaceholder}[1]{\vspace{0.2cm}\begin{center}\fbox{\begin{minipage}{0.82\textwidth}\centering\textcolor{blue!70!black}{\textbf{【插图位置】}}\\[0.2cm]#1\end{minipage}}\end{center}\vspace{0.2cm}}
\title{\textcolor{blue!70!black}{\Huge\textbf{生物2 · 生物体的结构层次}}}
\author{}\date{}
\begin{document}\maketitle\thispagestyle{empty}
\section*{\textcolor{blue!70!black}{致同学}}
生物1从分子和细胞讲起。现在我们要"放大镜头"——看看细胞怎么组成组织、组织怎么组成器官、器官怎么组成系统、系统怎么组成一个个完整的生物体。

\newpage
\section{\textcolor{orange!80!black}{第一课\quad 结构层次总论}}

\subsection*{从细胞到个体}
多细胞生物体的层级结构：
\formula{\text{细胞} \rightarrow \text{组织} \rightarrow \text{器官} \rightarrow \text{系统} \rightarrow \text{个体}}

\textbf{动物四大组织}：上皮组织（覆盖保护）、肌肉组织（收缩运动）、神经组织（传导信息）、结缔组织（连接支持——血液、骨头、韧带都是结缔组织）。\\
\textbf{植物三大组织}：保护组织（表皮）、疏导组织（导管/筛管）、营养组织（薄壁细胞）。

\subsection*{植物体 vs 动物体的结构特点}
\textbf{植物}：由六大器官组成（根茎叶花果种子）——但没有像动物那样的"系统"概念。\\
\textbf{动物}：由多个系统协同工作（消化/循环/呼吸/泌尿/神经/内分泌/运动/生殖）。

\begin{practice}
1. 生物体的结构层次：\_\_\_\_\_$\rightarrow$\_\_\_\_\_$\rightarrow$\_\_\_\_\_$\rightarrow$\_\_\_\_\_$\rightarrow$个体。\\
2. 人体中血液属于\_\_\_\_\_组织。\\
3. 植物体的六大器官：\_\_\_\_\_、\_\_\_\_\_、\_\_\_\_\_、\_\_\_\_\_、\_\_\_\_\_、\_\_\_\_\_。\\
4. 动物体具有而植物体不具有的结构层次是\_\_\_\_\_。\\
5. （判断）植物体没有"系统"这一结构层次。（\quad）\\
6. 保护组织属于\_\_\_\_\_（动物/植物）的组织类型。\\
7. 神经组织的主要功能是\_\_\_\_\_。\\
8. 结缔组织包括血液、骨头、\_\_\_\_\_等。
\end{practice}

\newpage
\section{\textcolor{orange!80!black}{第二课\quad 植物体的结构与功能}}

\subsection*{六器官分工}
\textbf{根}：吸收水分和无机盐、固定植物。根尖有根毛（增大吸收面积）。\\
\textbf{茎}：运输和支持。木质部（导管——运输水向上）、韧皮部（筛管——运输有机物向下）。\\
\textbf{叶}：光合作用的主要场所。气孔控制气体交换。\\
\textbf{花}：繁殖器官——雄蕊（花药产生花粉）、雌蕊（子房含胚珠）。\\
\textbf{果实}：保护种子、帮助种子传播。\\
\textbf{种子}：休眠的"小植物"——胚芽胚根子叶。

\subsection*{蒸腾作用}
植物体内的水通过叶片气孔以水蒸气形式散失——这叫\textbf{蒸腾作用}。它产生了一股向上的拉力——把水和矿物从根部一直吸到叶片。"蒸腾拉力"是植物运输水分的核心动力。叶片正面的气孔密度远小于背面——所以植物"脸上"水分蒸发起防晒作用。

\begin{practice}
1. 运输水分和无机盐的结构是\_\_\_\_\_（导管/筛管）；运输有机物的是\_\_\_\_\_。\\
2. （判断）蒸腾作用产生的拉力是植物从根部吸水上来的主要动力。（\quad）\\
3. 根尖有大量\_\_\_\_\_来增大吸收面积。\\
4. 一朵花的雌蕊中含有一或多个\_\_\_\_\_。\\
5. 叶片控制气体交换的结构是\_\_\_\_\_。\\
6. （判断）筛管运输水分和无机盐。（\quad）\\
7. 果实的主要功能是\_\_\_\_\_。\\
8. 种子中的胚芽将来发育成\_\_\_\_\_。
\end{practice}

\newpage
\section{\textcolor{orange!80!black}{第三课\quad 植物的主要生理}}

\subsection*{三大生理活动}
光合作用（在生物1已学，此处仅复习要点）：CO$_2$+H$_2$O $\xrightarrow{\text{光}}$ 葡萄糖+O$_2$。场所：叶绿体。\\
呼吸作用（复习）：葡萄糖+O$_2$ $\rightarrow$ CO$_2$+H$_2$O+能量。

\textbf{光合 vs 呼吸的关键对比：}
\begin{itemize}
  \item 光合：白天进行，\textbf{合成}有机物、\textbf{储存}能量。
  \item 呼吸：全天进行，\textbf{分解}有机物、\textbf{释放}能量（为生命活动供能）。
\end{itemize}

光合 $>$ 呼吸时——植物积累有机物，生长。光合 $=$ 呼吸时——净生长为零（光补偿点）。光合 $<$ 呼吸时——消耗储存，变瘦。

\subsection*{向性运动}
植物不能走路——但它们能\textbf{向特定方向生长}：\\
\textbf{向光性}：茎朝光生长（生长素在背光侧分布多→背光侧长得快→茎弯向光）。\\
\textbf{向地性}：根朝下长（向重力方向）、茎朝上长（背重力方向）。

\begin{thinkbox}
为什么窗台上的植物弯向窗外？因为窗外光照强——茎尖的生长素在背光一侧（室内）浓度更高——长得更快——茎就被"推"向了窗户方向。植物不是"主动转过去"的，而是"一侧长得快→弯了过去"。
\end{thinkbox}

\begin{practice}
1. 光合作用储存能量，呼吸作用\_\_\_\_\_能量。\\
2. （判断）植物白天只进行光合作用，晚上只进行呼吸作用。（\quad）\\
3. 植物向光性是因为背光侧生长素浓度\_\_\_\_\_（高/低），细胞长得\_\_\_\_\_（快/慢）。\\
4. 光补偿点是指光合\_\_\_\_\_呼吸。\\
5. 植物根向地生长，这属于\_\_\_\_\_性。\\
6. 植物进行光合作用的主要器官是\_\_\_\_\_。\\
7. 呼吸作用\_\_\_\_\_（全天/只在白天）进行。\\
8. 向光性中生长素在背光侧分布\_\_\_\_\_（多/少）。
\end{practice}

\newpage
\section{\textcolor{orange!80!black}{第四课\quad 动物的主要类群}}

\subsection*{无脊椎动物}
\begin{itemize}
  \item \textbf{腔肠动物}：水螅、水母。有口无肛门，辐射对称。
  \item \textbf{扁形动物}：涡虫、绦虫。身体扁平，有口无肛门。
  \item \textbf{线形动物}：蛔虫。圆柱形，有口有肛门。
  \item \textbf{环节动物}：蚯蚓。身体分节，有体腔。
  \item \textbf{软体动物}：蜗牛、河蚌、乌贼。身体柔软，常被贝壳。
  \item \textbf{节肢动物}：昆虫（蝗虫/蝴蝶）、蛛形纲（蜘蛛）、甲壳纲（虾蟹）、多足纲（蜈蚣）。身体分节、外骨骼（必须蜕皮才能长大）、有分节的附肢。是种类最多的动物类群。
\end{itemize}

\subsection*{脊椎动物}

\textbf{鱼类}：用鳃呼吸，用鳍游泳，体表有鳞片。体温随外界变化（变温动物）。\\
\textbf{两栖类}：幼体水中用鳃、成体用肺+皮肤辅助呼吸。青蛙、蟾蜍。变温。\\
\textbf{爬行类}：体表有角质鳞片或甲、用肺呼吸、陆地产卵（羊膜卵——脱离水体限制）。真正征服了陆地的脊椎动物。变温。\\
\textbf{鸟类}：体表覆羽、前肢成翼、有喙无齿、长骨中空、体温恒定（恒温）。\\
\textbf{哺乳类}：胎生哺乳（幼体在母体内发育、出生后用乳汁喂养）、体表被毛、恒温。人类的远房亲戚。

\subsection*{分类核心特征速记}
\begin{itemize}
  \item 鱼——鳃 + 鳍
  \item 两栖——水中幼体+陆地成体（变态）
  \item 爬行——羊膜卵 + 鳞片
  \item 鸟——羽毛 + 恒温 + 前肢成翼
  \item 哺乳——胎生哺乳 + 体毛
\end{itemize}

\begin{practice}
1. 种类最多的动物类群是\_\_\_\_\_动物。\\
2. （判断）两栖动物只能在水中生活。（\quad）\\
3. 真正摆脱了对水体依赖的脊椎动物是\_\_\_\_\_动物。\\
4. 哺乳动物最显著的特征是\_\_\_\_\_和\_\_\_\_\_。\\
5. 鱼类用\_\_\_\_\_呼吸，用\_\_\_\_\_游泳。\\
6. （判断）蝙蝠会飞，所以属于鸟类。（\quad）\\
7. 两栖动物的幼体生活在\_\_\_\_\_中。\\
8. 爬行动物体表覆盖\_\_\_\_\_。
\end{practice}

\newpage
\section{\textcolor{orange!80!black}{第五课\quad 人体八大系统（一）}}

\subsection*{消化系统}
从口到肛门一条"加工线"（已在本系列自然故事中详细讲过——此处复习）：\\
口腔（唾液淀粉酶）$\rightarrow$食道（蠕动）$\rightarrow$胃（胃酸+蛋白酶）$\rightarrow$小肠（消化吸收主要场所——绒毛增大表面积）$\rightarrow$大肠（回收水分）。

\subsection*{呼吸系统}
鼻$\rightarrow$咽$\rightarrow$喉$\rightarrow$气管$\rightarrow$支气管$\rightarrow$肺（肺泡——气体交换场所）。\\
肺泡壁极薄（一层细胞）、肺泡周围密布毛细血管——气体在此交换：O$_2$进血、CO$_2$出血。

\subsection*{循环系统}
\textbf{心脏}：四个腔室（左心房/左心室/右心房/右心室）。左心室壁最厚（把血泵到全身）。\\
\textbf{体循环}：左心室$\rightarrow$全身$\rightarrow$右心房（携带氧气给组织，带回CO$_2$）。\\
\textbf{肺循环}：右心室$\rightarrow$肺$\rightarrow$左心房（排出CO$_2$、获取O$_2$）。

\textbf{血液成分}：血浆（运载血细胞+营养物质）、红细胞（含血红蛋白——运输O$_2$）、白细胞（免疫）、血小板（凝血）。

\begin{figplaceholder}
此处插入心脏四腔与体循环/肺循环简图（见 figure/requirements/bio2-ch1.txt 插图1）
\end{figplaceholder}

\begin{practice}
1. 消化和吸收的主要场所是\_\_\_\_\_。\\
2. 气体交换的场所是\_\_\_\_\_。\\
3. （判断）血液呈红色是因为白细胞含有血红蛋白。（\quad）\\
4. 心脏四个腔中壁最厚的是\_\_\_\_\_。\\
5. 消化系统中吸收水分的主要场所是\_\_\_\_\_。\\
6. 肺泡壁由\_\_\_\_\_层细胞构成。\\
7. 红细胞的功能是运输\_\_\_\_\_。\\
8. 血小板的功能是\_\_\_\_\_。
\end{practice}

\newpage
\section{\textcolor{orange!80!black}{第六课\quad 人体八大系统（二）}}

\subsection*{泌尿系统}
肾脏（滤出尿液）$\rightarrow$输尿管$\rightarrow$膀胱（储存）$\rightarrow$尿道。\\
肾单位是肾脏的基本功能单位——包括肾小球（滤过）、肾小管（重吸收——把有用物质从原尿中"捡"回身体，留下废物形成尿液）。

\subsection*{神经系统}
\textbf{中枢神经}：脑 + 脊髓。\textbf{外周神经}：脑神经 + 脊神经。\\
反射弧（神经活动的"最小单元"）：感受器$\rightarrow$传入神经$\rightarrow$神经中枢$\rightarrow$传出神经$\rightarrow$效应器（肌肉或腺体）。

\subsection*{内分泌系统}
通过\textbf{激素}调节人体活动：\\
\textbf{甲状腺激素}：调节生长发育和代谢。缺碘$\rightarrow$甲状腺肿（大脖子病）。\\
\textbf{胰岛素}：降低血糖。分泌不足$\rightarrow$糖尿病。\\
\textbf{性激素}：促进性器官发育和第二性征（青春期变化）。

\subsection*{运动系统}
骨骼 + 关节 + 肌肉（骨骼肌附着在骨上——收缩时牵动骨运动）。肌肉只能\textbf{拉}骨不能\textbf{推}骨——所以肌肉多成对工作（如肱二头肌收缩→屈肘；肱三头肌收缩→伸肘）。

\subsection*{生殖系统}
男性的睾丸产生精子；女性的卵巢产生卵子。受精在输卵管内发生$\rightarrow$受精卵在子宫壁上着床$\rightarrow$发育成胚胎$\rightarrow$约38周后分娩。

\begin{practice}
1. 尿的形成包括\_\_\_\_\_和\_\_\_\_\_两个阶段。\\
2. 反射弧的五部分：\_\_\_\_\_$\rightarrow$\_\_\_\_\_$\rightarrow$\_\_\_\_\_$\rightarrow$\_\_\_\_\_$\rightarrow$\_\_\_\_\_。\\
3. 胰岛素分泌不足会导致\_\_\_\_\_病。\\
4. （判断）肌肉收缩时只能拉骨不能推骨。（\quad）\\
5. 肾单位包括\_\_\_\_\_和\_\_\_\_\_。\\
6. 反射弧中效应器是指\_\_\_\_\_（肌肉或腺体/神经）。\\
7. 胰岛素的作用是\_\_\_\_\_（升高/降低）血糖。\\
8. 受精作用发生在\_\_\_\_\_中。
\end{practice}

\newpage
\section{\textcolor{orange!80!black}{第七课\quad 动物的行为}}

\subsection*{先天性行为 vs 学习行为}
\textbf{先天性行为（本能）}：生来就会。蜘蛛织网、蜜蜂跳舞、小鸟求偶——不需要学习。\\
\textbf{学习行为}：通过经验和学习获得。鹦鹉学舌、小狗算算术、人类学骑车——需要练习。

学习行为需要大脑皮层的参与（高等动物学习能力强）。

\subsection*{社会行为}
许多动物营社会生活（蜜蜂、蚂蚁、狼群、猴群）——有明确分工、有等级或"语言"。蜜蜂通过\textbf{圆圈舞}和\textbf{8字舞}传递蜜源的方向和距离信息——这是动物中的"语言"。

\begin{practice}
1. "蜘蛛织网"是\_\_\_\_\_性行为，"鹦鹉学舌"是\_\_\_\_\_性行为。\\
2. 学习行为的优势是可以\_\_\_\_\_环境的变化。\\
3. （判断）动物越高等，先天性行为所占比例越大。（\quad）\\
4. 蜜蜂通过\_\_\_\_\_舞和\_\_\_\_\_舞传递蜜源信息。\\
5. 动物社会行为的主要特征是\_\_\_\_\_。\\
6. （判断）学习行为不需要大脑皮层的参与。（\quad）\\
7. 蜘蛛织网属于\_\_\_\_\_性行为。
\end{practice}

\newpage
\section{\textcolor{orange!80!black}{第八课\quad 微生物的世界}}

\subsection*{细菌——最简单的细胞生物}
单细胞原核生物（没有真正的细胞核）。分裂繁殖（一个变两个，约20分钟一代）。按形态分：球菌、杆菌、螺旋菌。依靠\textbf{细胞壁}保护、\textbf{鞭毛}运动。在自然界中广泛存在——土壤中每克有上亿个细菌。

\subsection*{真菌}
酵母菌（单细胞真菌——用于酿酒和发酵）；霉菌（青霉菌产生青霉素——第一种抗生素）；蘑菇（大型真菌，多细胞）。真菌是真核生物，有真正的细胞核。

\subsection*{病毒}
极小——比细菌还小得多（纳米级）。结构简单——蛋白质外壳 + 内部遗传物质（DNA或RNA）。不能独立繁殖——只能侵入活细胞，利用宿主细胞的"机器"来复制自己。

\textbf{病毒引起的常见疾病}：流感、新冠、乙肝、艾滋病。治疗病毒远比治疗细菌难——因为病毒藏在你的细胞内部，抗生素对病毒无效。

\begin{thinkbox}
为什么普通感冒吃抗生素没用？因为感冒通常由病毒引起——抗生素只对\textbf{细菌}有效。滥用抗生素会筛选出"超级细菌"（耐药菌）。所以医生反复强调：不要自己随便吃抗生素。
\end{thinkbox}

\begin{practice}
1. 细菌是\_\_\_\_\_核生物；真菌是\_\_\_\_\_核生物。\\
2. 抗生素对\_\_\_\_\_有效，对\_\_\_\_\_无效。\\
3. 病毒由\_\_\_\_\_和\_\_\_\_\_两部分组成。\\
4. （判断）所有微生物都对人体有害。（\quad）\\
5. 细菌的繁殖方式是\_\_\_\_\_。\\
6. （判断）病毒有完整的细胞结构。（\quad）\\
7. 青霉菌能产生抗生素\_\_\_\_\_。\\
8. 酵母菌属于\_\_\_\_\_（细菌/真菌）。
\end{practice}

\vspace{0.5cm}
\begin{center}\rule{0.7\textwidth}{1pt}\vspace{0.4cm}
{\large\textcolor{blue!70!black}{\textbf{—— 生物2·生物体的结构层次 完 ——}}}
\end{center}

\newpage
\section*{\textcolor{defblue}{参考答案}}

\subsection*{第一课}
1. 细胞$\rightarrow$组织$\rightarrow$器官$\rightarrow$系统$\rightarrow$个体\\
2. 结缔\\
3. 根、茎、叶、花、果实、种子\\
4. 系统\\
5. （√）\\
6. 植物\\
7. 传导（传递）信息\\
8. 韧带

\subsection*{第二课}
1. 导管、筛管\\
2. （√）\\
3. 根毛\\
4. 胚珠\\
5. 气孔\\
6. （×）筛管运输有机物（向下）\\
7. 保护种子、帮助种子传播\\
8. 茎和叶

\subsection*{第三课}
1. 释放\\
2. （×）白天光合呼吸同时进行，晚上只有呼吸\\
3. 高、快\\
4. 等于\\
5. 向地\\
6. 叶\\
7. 全天\\
8. 多

\subsection*{第四课}
1. 节肢\\
2. （×）成体可在陆地生活，但繁殖离不开水\\
3. 爬行\\
4. 胎生、哺乳\\
5. 鳃、鳍\\
6. （×）蝙蝠是哺乳动物\\
7. 水\\
8. 角质鳞片或甲

\subsection*{第五课}
1. 小肠\\
2. 肺（肺泡）\\
3. （×）血红蛋白在红细胞中，使血液呈红色\\
4. 左心室\\
5. 大肠\\
6. 一\\
7. 氧气\\
8. 凝血（止血）

\subsection*{第六课}
1. 肾小球的滤过、肾小管的重吸收\\
2. 感受器$\rightarrow$传入神经$\rightarrow$神经中枢$\rightarrow$传出神经$\rightarrow$效应器\\
3. 糖尿\\
4. （√）\\
5. 肾小球、肾小管\\
6. 肌肉或腺体\\
7. 降低\\
8. 输卵管

\subsection*{第七课}
1. 先天、学习\\
2. 适应\\
3. （×）越高等学习行为比例越大\\
4. 圆圈（舞）、8字（舞）\\
5. 有明确分工、有等级或"语言"\\
6. （×）学习行为需要大脑皮层参与\\
7. 先天

\subsection*{第八课}
1. 原、真\\
2. 细菌、病毒\\
3. 蛋白质外壳、内部遗传物质（DNA或RNA）\\
4. （×）很多微生物有益\\
5. 分裂（二分裂）\\
6. （×）病毒没有细胞结构\\
7. 青霉素\\
8. 真菌

\end{document}
